\section{SAR image time series model}
    \label{section:data_model}
    \noindent
    For a given stack of $l = p+1$ \acs{SAR} images, we consider a local homogeneous spatial neighborhood of size $n$ denoted $\{\mathbf{\widetilde{x}}^i\}_{i=1}^n$ where $\mathbf{\widetilde{x}}^i \in \mathbb{C}^{l}$ for all $i \in [\! [1,n] \!]$, i.e., 
    \begin{equation}
        \label{eq:data}
            \mathbf{\widetilde{x}}^i = [\underbrace{x_1^i, \dots, x_p^i}_{\mathbf{x}^i}, x_{l}^i ]^T \in \mathbb{C}^{l},
    \end{equation}
    \noindent
    $\mathbf{x}^i \in \mathbb{C}^{p}$ indicates the multivariate pixel of the previous images (Fig. \ref{fig:datacube}). %\footnote{\footnotesize \textcolor{red}{The tilde represents all the data, and when absent, it refers to the past data, whereas the hat symbol indicates that the parameter used is an estimate.}} (Fig. \ref{fig:datacube}). 
    \vspace{-0.5cm}
    \begin{figure}[hbt]
        \centering
        

\tikzset{every picture/.style={line width=0.75pt}} %set default line width to 0.75pt        

\begin{tikzpicture}[x=0.35pt,y=0.35pt,yscale=-1,xscale=1]
%uncomment if require: \path (0,625); %set diagram left start at 0, and has height of 625



%Shape: Grid [id:dp3363536071719442] 
\draw  [draw opacity=0][fill={rgb, 255:red, 255; green, 255; blue, 255 }  ,fill opacity=1 ] (97.48,336.33) -- (256.41,336.33) -- (256.41,495.8) -- (97.48,495.8) -- cycle ; \draw   (97.48,336.33) -- (97.48,495.8)(111.9,336.33) -- (111.9,495.8)(126.31,336.33) -- (126.31,495.8)(140.73,336.33) -- (140.73,495.8)(155.14,336.33) -- (155.14,495.8)(169.56,336.33) -- (169.56,495.8)(183.97,336.33) -- (183.97,495.8)(198.39,336.33) -- (198.39,495.8)(212.8,336.33) -- (212.8,495.8)(227.22,336.33) -- (227.22,495.8)(241.63,336.33) -- (241.63,495.8)(256.05,336.33) -- (256.05,495.8) ; \draw   (97.48,336.33) -- (256.41,336.33)(97.48,350.74) -- (256.41,350.74)(97.48,365.16) -- (256.41,365.16)(97.48,379.57) -- (256.41,379.57)(97.48,393.99) -- (256.41,393.99)(97.48,408.4) -- (256.41,408.4)(97.48,422.82) -- (256.41,422.82)(97.48,437.23) -- (256.41,437.23)(97.48,451.65) -- (256.41,451.65)(97.48,466.06) -- (256.41,466.06)(97.48,480.48) -- (256.41,480.48)(97.48,494.89) -- (256.41,494.89) ; \draw    ;
%Straight Lines [id:da8990331411349495] 
\draw    (100.48,512.7) -- (254.8,512.46) ;
\draw [shift={(257.8,512.46)}, rotate = 179.91] [fill={rgb, 255:red, 0; green, 0; blue, 0 }  ][line width=0.08]  [draw opacity=0] (10.72,-5.15) -- (0,0) -- (10.72,5.15) -- (7.12,0) -- cycle    ;
\draw [shift={(97.48,512.7)}, rotate = 359.91] [fill={rgb, 255:red, 0; green, 0; blue, 0 }  ][line width=0.08]  [draw opacity=0] (8.93,-4.29) -- (0,0) -- (8.93,4.29) -- cycle    ;
%Straight Lines [id:da6084690298752402] 
\draw    (84.2,340.6) -- (84.2,493.6) ;
\draw [shift={(84.2,496.6)}, rotate = 270] [fill={rgb, 255:red, 0; green, 0; blue, 0 }  ][line width=0.08]  [draw opacity=0] (10.72,-5.15) -- (0,0) -- (10.72,5.15) -- (7.12,0) -- cycle    ;
\draw [shift={(84.2,337.6)}, rotate = 90] [fill={rgb, 255:red, 0; green, 0; blue, 0 }  ][line width=0.08]  [draw opacity=0] (8.93,-4.29) -- (0,0) -- (8.93,4.29) -- cycle    ;
%Shape: Grid [id:dp019122537703205822] 
\draw  [draw opacity=0][fill={rgb, 255:red, 128; green, 128; blue, 128 }  ,fill opacity=1 ] (126.21,365.16) -- (170.39,365.16) -- (170.39,409.28) -- (126.21,409.28) -- cycle ; \draw   (126.21,365.16) -- (126.21,409.28)(140.62,365.16) -- (140.62,409.28)(155.04,365.16) -- (155.04,409.28)(169.45,365.16) -- (169.45,409.28) ; \draw   (126.21,365.16) -- (170.39,365.16)(126.21,379.57) -- (170.39,379.57)(126.21,393.99) -- (170.39,393.99)(126.21,408.4) -- (170.39,408.4) ; \draw    ;

%Shape: Grid [id:dp9622630659884281] 
\draw  [draw opacity=0][fill={rgb, 255:red, 255; green, 255; blue, 255 }  ,fill opacity=1 ] (391.4,32.84) -- (550.33,32.84) -- (550.33,192.31) -- (391.4,192.31) -- cycle ; \draw   (391.4,32.84) -- (391.4,192.31)(405.82,32.84) -- (405.82,192.31)(420.23,32.84) -- (420.23,192.31)(434.65,32.84) -- (434.65,192.31)(449.06,32.84) -- (449.06,192.31)(463.48,32.84) -- (463.48,192.31)(477.89,32.84) -- (477.89,192.31)(492.31,32.84) -- (492.31,192.31)(506.72,32.84) -- (506.72,192.31)(521.14,32.84) -- (521.14,192.31)(535.55,32.84) -- (535.55,192.31)(549.97,32.84) -- (549.97,192.31) ; \draw   (391.4,32.84) -- (550.33,32.84)(391.4,47.25) -- (550.33,47.25)(391.4,61.67) -- (550.33,61.67)(391.4,76.08) -- (550.33,76.08)(391.4,90.5) -- (550.33,90.5)(391.4,104.91) -- (550.33,104.91)(391.4,119.33) -- (550.33,119.33)(391.4,133.74) -- (550.33,133.74)(391.4,148.16) -- (550.33,148.16)(391.4,162.57) -- (550.33,162.57)(391.4,176.99) -- (550.33,176.99)(391.4,191.4) -- (550.33,191.4) ; \draw    ;
%Shape: Cube [id:dp5751813388717848] 
\draw  [fill={rgb, 255:red, 128; green, 128; blue, 128 }  ,fill opacity=1 ] (319.98,161.76) -- (420.31,61.43) -- (463.48,61.43) -- (463.48,104.91) -- (363.15,205.25) -- (319.98,205.25) -- cycle ; \draw   (463.48,61.43) -- (363.15,161.76) -- (319.98,161.76) ; \draw   (363.15,161.76) -- (363.15,205.25) ;
%Shape: Grid [id:dp14911386097810908] 
\draw  [draw opacity=0][fill={rgb, 255:red, 255; green, 255; blue, 255 }  ,fill opacity=1 ][line width=0.75]  (290.5,133.74) -- (449.42,133.74) -- (449.42,293.21) -- (290.5,293.21) -- cycle ; \draw  [line width=0.75]  (290.5,133.74) -- (290.5,293.21)(304.91,133.74) -- (304.91,293.21)(319.33,133.74) -- (319.33,293.21)(333.74,133.74) -- (333.74,293.21)(348.16,133.74) -- (348.16,293.21)(362.57,133.74) -- (362.57,293.21)(376.99,133.74) -- (376.99,293.21)(391.4,133.74) -- (391.4,293.21)(405.82,133.74) -- (405.82,293.21)(420.23,133.74) -- (420.23,293.21)(434.65,133.74) -- (434.65,293.21)(449.06,133.74) -- (449.06,293.21) ; \draw  [line width=0.75]  (290.5,133.74) -- (449.42,133.74)(290.5,148.16) -- (449.42,148.16)(290.5,162.57) -- (449.42,162.57)(290.5,176.99) -- (449.42,176.99)(290.5,191.4) -- (449.42,191.4)(290.5,205.82) -- (449.42,205.82)(290.5,220.23) -- (449.42,220.23)(290.5,234.65) -- (449.42,234.65)(290.5,249.06) -- (449.42,249.06)(290.5,263.48) -- (449.42,263.48)(290.5,277.89) -- (449.42,277.89)(290.5,292.31) -- (449.42,292.31) ; \draw  [line width=0.75]   ;
%Straight Lines [id:da9954186422316147] 
\draw  [dash pattern={on 0.84pt off 2.51pt}]  (450.05,291.13) -- (550.58,190.22) ;
%Straight Lines [id:da125892397375575] 
\draw  [dash pattern={on 0.84pt off 2.51pt}]  (292.08,132.57) -- (392.61,31.66) ;
%Straight Lines [id:da5438545149879253] 
\draw  [dash pattern={on 0.84pt off 2.51pt}]  (450.05,132.57) -- (550.58,31.66) ;
%Shape: Cube [id:dp4454804477484464] 
\draw  [fill={rgb, 255:red, 128; green, 128; blue, 128 }  ,fill opacity=0.46 ] (319.9,162) -- (420.23,61.67) -- (463.4,61.67) -- (463.4,105.15) -- (363.07,205.48) -- (319.9,205.48) -- cycle ; \draw   (463.4,61.67) -- (363.07,162) -- (319.9,162) ; \draw   (363.07,162) -- (363.07,205.48) ;
%Straight Lines [id:da8853614458757708] 
\draw    (274.23,135.4) -- (274,291) ;
\draw [shift={(274,294)}, rotate = 270.08] [fill={rgb, 255:red, 0; green, 0; blue, 0 }  ][line width=0.08]  [draw opacity=0] (10.72,-5.15) -- (0,0) -- (10.72,5.15) -- (7.12,0) -- cycle    ;
\draw [shift={(274.24,132.4)}, rotate = 90.08] [fill={rgb, 255:red, 0; green, 0; blue, 0 }  ][line width=0.08]  [draw opacity=0] (8.93,-4.29) -- (0,0) -- (8.93,4.29) -- cycle    ;
%Straight Lines [id:da6253357160499429] 
\draw    (287.45,311.07) -- (441.76,310.83) ;
\draw [shift={(444.76,310.83)}, rotate = 179.91] [fill={rgb, 255:red, 0; green, 0; blue, 0 }  ][line width=0.08]  [draw opacity=0] (10.72,-5.15) -- (0,0) -- (10.72,5.15) -- (7.12,0) -- cycle    ;
\draw [shift={(284.45,311.07)}, rotate = 359.91] [fill={rgb, 255:red, 0; green, 0; blue, 0 }  ][line width=0.08]  [draw opacity=0] (8.93,-4.29) -- (0,0) -- (8.93,4.29) -- cycle    ;
%Shape: Grid [id:dp11960772992405277] 
\draw  [draw opacity=0][fill={rgb, 255:red, 128; green, 128; blue, 128 }  ,fill opacity=1 ] (319.9,162.24) -- (364.09,162.24) -- (364.09,206.36) -- (319.9,206.36) -- cycle ; \draw   (319.9,162.24) -- (319.9,206.36)(334.32,162.24) -- (334.32,206.36)(348.73,162.24) -- (348.73,206.36)(363.15,162.24) -- (363.15,206.36) ; \draw   (319.9,162.24) -- (364.09,162.24)(319.9,176.65) -- (364.09,176.65)(319.9,191.07) -- (364.09,191.07)(319.9,205.48) -- (364.09,205.48) ; \draw    ;
%Shape: Grid [id:dp9370576522410461] 
\draw  [draw opacity=0][fill={rgb, 255:red, 128; green, 128; blue, 128 }  ,fill opacity=0.27 ] (420.23,61.43) -- (464.42,61.43) -- (464.42,105.55) -- (420.23,105.55) -- cycle ; \draw  [color={rgb, 255:red, 0; green, 0; blue, 0 }  ,draw opacity=0.35 ] (420.23,61.43) -- (420.23,105.55)(434.65,61.43) -- (434.65,105.55)(449.06,61.43) -- (449.06,105.55)(463.48,61.43) -- (463.48,105.55) ; \draw  [color={rgb, 255:red, 0; green, 0; blue, 0 }  ,draw opacity=0.35 ] (420.23,61.43) -- (464.42,61.43)(420.23,75.85) -- (464.42,75.85)(420.23,90.26) -- (464.42,90.26)(420.23,104.68) -- (464.42,104.68) ; \draw  [color={rgb, 255:red, 0; green, 0; blue, 0 }  ,draw opacity=0.35 ]  ;

%Straight Lines [id:da6271920605962396] 
\draw    (234.11,597.87) -- (597.3,231.1) ;
\draw [shift={(232,600)}, rotate = 314.72] [fill={rgb, 255:red, 0; green, 0; blue, 0 }  ][line width=0.08]  [draw opacity=0] (8.93,-4.29) -- (0,0) -- (8.93,4.29) -- cycle    ;
%Straight Lines [id:da1054085848226396] 
\draw    (292.55,530.67) -- (302.49,543.38) ;


%Straight Lines [id:da21811336543584559] 
\draw    (588.47,230.59) -- (598.41,243.3) ;

%Straight Lines [id:da7877364989769449] 
\draw    (485.73,332.61) -- (495.67,345.31) ;

%Curve Lines [id:da1317888236393614] 
\draw    (389,113) .. controls (359.15,48.32) and (317.42,-22.29) .. (287.45,49.9) ;
\draw [shift={(287,51)}, rotate = 292.07] [color={rgb, 255:red, 0; green, 0; blue, 0 }  ][line width=0.75]    (10.93,-4.9) .. controls (6.95,-2.3) and (3.31,-0.67) .. (0,0) .. controls (3.31,0.67) and (6.95,2.3) .. (10.93,4.9)   ;
%Curve Lines [id:da9274466791250591] 
\draw    (144.73,386.57) .. controls (74.71,301.86) and (52.44,280.43) .. (107.31,260.6) ;
\draw [shift={(109,260)}, rotate = 160.67] [color={rgb, 255:red, 0; green, 0; blue, 0 }  ][line width=0.75]    (10.93,-3.29) .. controls (6.95,-1.4) and (3.31,-0.3) .. (0,0) .. controls (3.31,0.3) and (6.95,1.4) .. (10.93,3.29)   ;


% Text Node
\draw (385,340) node  [font=\small] [align=left] {\begin{minipage}[lt]{41.32pt}\setlength\topsep{0pt}
{\normalsize Spatial}
\end{minipage}};
% Text Node
\draw (256.44,190) node  [font=\small,rotate=-270.23] [align=left] {\begin{minipage}[lt]{46.1pt}\setlength\topsep{0pt}
{\normalsize Spatial}
\end{minipage}};
% Text Node
\draw (380,1) node [anchor=north west][inner sep=0.75pt]   [align=left] {{\normalsize Previous images}};
% Text Node
\draw (280,555) node [anchor=north west][inner sep=0.75pt]    {$t_{p+1}$};
% Text Node
\draw (427.63,450) node  [font=\small,rotate=-315.27,xslant=0] [align=left] {\begin{minipage}[lt]{77.06pt}\setlength\topsep{0pt}
{\normalsize Temporal axis}\\
\end{minipage}};
% Text Node
\draw (605,250) node [anchor=north west][inner sep=0.75pt]    {$t_{0}$};
% Text Node
\draw (500,360) node [anchor=north west][inner sep=0.75pt]    {$t_{p}$};
% Text Node
\draw (195,540) node  [font=\small] [align=left] {\begin{minipage}[lt]{41.32pt}\setlength\topsep{0pt}
{\normalsize Spatial}
\end{minipage}};
% Text Node
\draw (60,400) node  [font=\small,rotate=-270.23] [align=left] {\begin{minipage}[lt]{41.47pt}\setlength\topsep{0pt}
{\normalsize Spatial}
\end{minipage}};
% Text Node
\draw (100,300) node [anchor=north west][inner sep=0.75pt]   [align=left] {{\normalsize New \ image}};
% Text Node
\draw (120,255) node [anchor=north west][inner sep=0.75pt]    {$\{x_{l}^{i} \}_{i=1}^{n}$};
% Text Node
\draw (100,226) node [anchor=north west][inner sep=0.75pt]   [align=left] {local patch};
% Text Node
\draw (200,85) node [anchor=north west][inner sep=0.75pt]    {$\{\mathbf{x}^{i} \}_{i=1}^{n}$};
% Text Node
\draw (195,57) node [anchor=north west][inner sep=0.75pt]   [align=left] {local patch};

\end{tikzpicture}

        %\vspace{-10pt}
        \caption{\small The representation of \acs{SAR} data incorporates both previous and recently acquired images. The local neighborhood of size $n$ is denoted by gray pixels (sliding window).}
        \label{fig:datacube}
    \end{figure}
    \vspace{-0.2cm}
    A common assumption consists in considering that $\{\mathbf{\widetilde{x}}^i\}_{i=1}^n$ is a set of vectors \ac{i.i.d} following a zero mean \acf{CCG} distribution \citep{bamler1998synthetic}, i.e., $\mathbf{\widetilde{x}}\sim\mathcal{CN}(0, \mathbf{\widetilde{\Sigma}})$. Most of \acs{PL} algorithms in the literature are based on this assumption \citep{ferretti2001permanent, ansari2018efficient, vu2022new, 9763551}, however \citep{mian2018new, vu2023robust} showed that this assumption does not always hold true
    %a Generalized Gaussian distribution provides a better fit for the Sentinel-1 image data 
    and addressed Scaled Gaussian model. For this scenario, the data follows a Gaussian distribution conditionally to a deterministic unknown scale $\tau_i$ \citep{ollila2012complex, maronna2019robust}, i.e,  $\mathbf{\widetilde{x}}^i \sim\mathcal{CN}(0, \tau_i \mathbf{\widetilde{\Sigma}})$.
    The real core of the covariance matrix is denoted by $\mathbf{\widetilde{\Psi}}$, and $\mathbf{\widetilde{w}}_{\mathbf{\theta}}$ represents the vector of the exponential of the complex phases ($\mathbf{\widetilde{w}}_{\mathbf{\theta}}=[e^{j \theta_0}, \dots,e^{j \theta_{l}}]$). Based on standard physical parameters in \acs{InSAR}, the covariance matrix respects the given structure \citep{guarnieri2008exploitation}
    \vspace{-0.4em}
    \begin{equation}
        \label{Equation:cov_mat_struc}
       \mathbf{\widetilde{\Sigma}} = \mathbf{\widetilde{\Psi}} \circ \mathbf{\widetilde{w}}_{\theta}\mathbf{\widetilde{w}}_{\theta}^H.
    \end{equation}
    %where the symbol $\circ$ characterizes the element-wise (Hadamard) multiplication, the exponent $^H$ corresponds to the transposed and complex conjugated operator (Hermitian). % $\mathbf{\widetilde{\Psi}}$ denotes the real core of the covariance matrix and $\mathbf{\widetilde{w}}_{\mathbf{\theta}}$ represents the vector of the exponential of the complex phases ($\mathbf{\widetilde{w}}_{\mathbf{\theta}}=[e^{j \theta_0}, \dots,e^{j \theta_{l}}]$).
    According to (\ref{Equation:cov_mat_struc}) and given the representation of the data in (\ref{eq:data}), the covariance matrix can be equivalently represented as 
    \begin{equation}
        \label{cov_mat_struc_block}
        \mathbf{\widetilde{\Sigma}} = \left( 
            \begin{array}{cc}
               \boldsymbol{\Sigma} &   w_{\theta_{l}}^* \mathrm{diag}(\mathbf{w}_{\theta})\boldsymbol{\gamma}^T  \\
                \boldsymbol{\gamma} \mathrm{diag}(\mathbf{w}_{\theta})^H w_{\theta_{l}}    & \gamma_{l} \\ 
            \end{array} \right),
    \end{equation}
    where $\mathbf{\Sigma}$ corresponds to the previously estimated covariance matrix between the previous \acs{SAR} images, $\boldsymbol{\gamma}$ denotes the coherence vector between the newly acquired data and the previous ones, $\gamma_{l}$ symbolizes the variance of the newly acquired data,  and $w_{\theta_{l}}$ is the exponential of the phase of the latest date %(with respect to the reference date) 
    and $ \mathbf{w}_{\theta}$ corresponds to the vector of the exponential of the complex phases differences of the past. %The tilde represents all the data, and when absent, it refers to the past data, whereas the hat symbol indicates that the parameter used is an estimate.
    %\textcolor{blue}{The exponent $^T$ (resp. $^*$) is  the transposed operator (resp. the conjugated operator).}